\qquad A partir da implementação e dos experimentos realizados com os
quatro algoritmos, foi possível observar, na prática, diferenças de
comportamento que confirmam a análise teórica de cada um.

\qquad O \textbf{Insertion Sort} e o \textbf{Bubble Sort} compartilham uma
característica importante: ambos são sensíveis ao grau de ordenação
da entrada, apresentando desempenho linear ($O(n)$) quando os dados já
estão ordenados (entrada crescente) e desempenho quadrático ($O(n^2)$)
nos casos médio e pior (entradas randômica e decrescente,
respectivamente). Entre os dois, o Bubble Sort mostrou-se
consistentemente mais lento que o Insertion Sort para os mesmos
tamanhos e tipos de entrada, apesar de ambos terem a mesma ordem de
grandeza de comparações no pior e no caso médio — diferença
explicada pelo maior custo de cada troca do Bubble Sort (três
atribuições) frente ao deslocamento único do Insertion Sort.

\qquad O \textbf{Selection Sort} se destacou por um comportamento distinto:
sua complexidade é $\Theta(n^2)$ em \textit{todos} os casos (melhor, médio e
pior), já que o número de comparações não depende da ordem dos dados
de entrada — apenas o número de trocas varia, mas essa quantidade é
sempre $O(n)$ e, portanto, irrelevante diante das $O(n^2)$ comparações.
Isso foi confirmado experimentalmente: os tempos medidos para as
entradas crescente, decrescente e randômica ficaram praticamente
idênticos entre si, para um mesmo tamanho de instância.

\qquad O \textbf{Shell Sort} apresentou, de longe, o melhor desempenho entre
os quatro algoritmos, para todos os tipos de entrada testados, mesmo
utilizando a sequência de intervalos mais simples (a original de
Shell). Embora essa sequência tenha pior caso teórico $\Theta(n^2)$, esse
pior caso exige uma disposição de dados especificamente construída
para isso, o que não ocorre nem mesmo com uma entrada totalmente
decrescente — por isso, tanto a entrada crescente quanto a
decrescente exibiram comportamento próximo de O(n log n) nos testes
realizados, enquanto a entrada randômica, embora sem uma fórmula
fechada simples de caso médio, também se manteve ordens de grandeza
mais rápida que os demais algoritmos para os mesmos tamanhos.

\qquad Em resumo, os experimentos evidenciam que, embora Insertion,
Selection e Bubble Sort compartilhem a mesma complexidade assintótica
de pior caso ($O(n^2)$), a escolha entre eles pode fazer diferença
prática relevante dependendo do grau de ordenação esperado dos dados:
o Insertion Sort é preferível quando há expectativa de entradas quase
ordenadas, o Selection Sort tem a vantagem de um número previsível e
mínimo de trocas (útil quando a operação de escrita é cara), e o
Bubble Sort, apesar de didático, não apresenta vantagem prática clara
sobre os demais nos cenários testados. O Shell Sort, por sua vez,
demonstra como uma modificação relativamente simples do Insertion
Sort (a introdução de saltos/gaps decrescentes) é suficiente para
alcançar um ganho de desempenho expressivo, sendo a escolha mais
adequada dentre os quatro para instâncias grandes, especialmente
quando não é garantida nenhuma característica sobre a ordenação
prévia dos dados.
